\section{支撑材料清单与完整源程序}

\subsection{支撑材料文件列表}

支撑材料压缩包采用如下目录结构；赛题提供的原始 Excel 不重复打包，复现时按 \texttt{README.md} 放入 \texttt{user\_data/} 即可。

\begin{longtable}{p{0.34\textwidth}p{0.60\textwidth}}
\toprule
路径 & 内容与用途 \\
\midrule
\endfirsthead
\toprule
路径 & 内容与用途 \\
\midrule
\endhead
\texttt{README.md} & 从空环境开始的安装、运行、输出和论文编译说明 \\
\texttt{requirements-lock.txt} & Python 3.12 的直接依赖精确版本 \\
\texttt{reproduce.py} & 四问、灵敏度、约束复核和图表的一键入口 \\
\texttt{code/} & 问题一至问题四、参数、工具、数据检查、灵敏度和约束复核的完整 Python 源码 \\
\texttt{figures/*.py}、\texttt{\_utils/} & 图表数据整理、统一绘图和 LaTeX 表格生成程序 \\
\texttt{paper/} & 论文 LaTeX 源文件、模板、字体、参考文献和已生成图表 \\
\texttt{output/} & 正式 CSV 及问题四可选比较输出的同内容副本 \\
\texttt{figures/problem\_*.json} & 四问完整机器结果、求解器状态、边界、审计和敏感性记录 \\
\path{user_data/worldcup_2022_openfootball.json} & 问题四真实赛程的机器读取快照 \\
\path{user_data/worldcup_2022_pre_tournament_elo.csv} & 问题四完整赛前 Elo 快照 \\
\texttt{actual\_schedule\_2022.csv} & 真实赛程逐场整理表、来源 URL 与访问日期 \\
\texttt{AI 工具使用详情.pdf} & 工具、用途、提示方式、采纳和核验记录 \\
\bottomrule
\end{longtable}

\subsection{完整结果附表}

以下附表由正式 JSON/CSV 直接生成，用于复核正文汇总值。

{
\footnotesize
\hfuzz=0.2pt
\setlength{\tabcolsep}{3pt}
\begin{longtable}{p{0.15\textwidth}p{0.16\textwidth}p{0.15\textwidth}p{0.18\textwidth}p{0.19\textwidth}}
\caption{问题一特征字典与时间血缘}\label{tab:q1-data-dictionary}\\
\toprule
特征 & 来源表 & 来源列 & 时间规则 & 拟合样本 \\
\midrule
\endfirsthead
\multicolumn{5}{c}{续表} \\
\toprule
特征 & 来源表 & 来源列 & 时间规则 & 拟合样本 \\
\midrule
\endhead
\midrule \multicolumn{5}{r}{续下页} \\
\endfoot
\bottomrule
\endlastfoot
competition & historical\_\allowbreak{}matches/\allowbreak{}teams/\allowbreak{}base\_\allowbreak{}predictions & competition & pregame\_\allowbreak{}available & outer\_\allowbreak{}fold\_\allowbreak{}train\_\allowbreak{}only;\allowbreak{} final=560 train ids \\
stage & historical\_\allowbreak{}matches/\allowbreak{}teams/\allowbreak{}base\_\allowbreak{}predictions & stage & pregame\_\allowbreak{}available & outer\_\allowbreak{}fold\_\allowbreak{}train\_\allowbreak{}only;\allowbreak{} final=560 train ids \\
timezone\_\allowbreak{}pair & historical\_\allowbreak{}matches/\allowbreak{}teams/\allowbreak{}base\_\allowbreak{}predictions & timezone\_\allowbreak{}pair & pregame\_\allowbreak{}available & outer\_\allowbreak{}fold\_\allowbreak{}train\_\allowbreak{}only;\allowbreak{} final=560 train ids \\
strength\_\allowbreak{}rank\_\allowbreak{}mean & historical\_\allowbreak{}matches/\allowbreak{}teams/\allowbreak{}base\_\allowbreak{}predictions & strength\_\allowbreak{}rank\_\allowbreak{}mean & pregame\_\allowbreak{}available & outer\_\allowbreak{}fold\_\allowbreak{}train\_\allowbreak{}only;\allowbreak{} final=560 train ids \\
strength\_\allowbreak{}rank\_\allowbreak{}gap & historical\_\allowbreak{}matches/\allowbreak{}teams/\allowbreak{}base\_\allowbreak{}predictions & strength\_\allowbreak{}rank\_\allowbreak{}gap & pregame\_\allowbreak{}available & outer\_\allowbreak{}fold\_\allowbreak{}train\_\allowbreak{}only;\allowbreak{} final=560 train ids \\
elo\_\allowbreak{}mean & historical\_\allowbreak{}matches/\allowbreak{}teams/\allowbreak{}base\_\allowbreak{}predictions & elo\_\allowbreak{}mean & pregame\_\allowbreak{}available & outer\_\allowbreak{}fold\_\allowbreak{}train\_\allowbreak{}only;\allowbreak{} final=560 train ids \\
elo\_\allowbreak{}gap & historical\_\allowbreak{}matches/\allowbreak{}teams/\allowbreak{}base\_\allowbreak{}predictions & elo\_\allowbreak{}gap & pregame\_\allowbreak{}available & outer\_\allowbreak{}fold\_\allowbreak{}train\_\allowbreak{}only;\allowbreak{} final=560 train ids \\
p\_\allowbreak{}a & historical\_\allowbreak{}matches/\allowbreak{}teams/\allowbreak{}base\_\allowbreak{}predictions & p\_\allowbreak{}a & pregame\_\allowbreak{}available & outer\_\allowbreak{}fold\_\allowbreak{}train\_\allowbreak{}only;\allowbreak{} final=560 train ids \\
p\_\allowbreak{}draw & historical\_\allowbreak{}matches/\allowbreak{}teams/\allowbreak{}base\_\allowbreak{}predictions & p\_\allowbreak{}draw & pregame\_\allowbreak{}available & outer\_\allowbreak{}fold\_\allowbreak{}train\_\allowbreak{}only;\allowbreak{} final=560 train ids \\
p\_\allowbreak{}b & historical\_\allowbreak{}matches/\allowbreak{}teams/\allowbreak{}base\_\allowbreak{}predictions & p\_\allowbreak{}b & pregame\_\allowbreak{}available & outer\_\allowbreak{}fold\_\allowbreak{}train\_\allowbreak{}only;\allowbreak{} final=560 train ids \\
odds\_\allowbreak{}entropy & historical\_\allowbreak{}matches/\allowbreak{}teams/\allowbreak{}base\_\allowbreak{}predictions & odds\_\allowbreak{}entropy & pregame\_\allowbreak{}available & outer\_\allowbreak{}fold\_\allowbreak{}train\_\allowbreak{}only;\allowbreak{} final=560 train ids \\
market\_\allowbreak{}value\_\allowbreak{}musd\_\allowbreak{}mean & historical\_\allowbreak{}matches/\allowbreak{}teams/\allowbreak{}base\_\allowbreak{}predictions & market\_\allowbreak{}value\_\allowbreak{}musd\_\allowbreak{}mean & pregame\_\allowbreak{}available & outer\_\allowbreak{}fold\_\allowbreak{}train\_\allowbreak{}only;\allowbreak{} final=560 train ids \\
market\_\allowbreak{}value\_\allowbreak{}musd\_\allowbreak{}gap & historical\_\allowbreak{}matches/\allowbreak{}teams/\allowbreak{}base\_\allowbreak{}predictions & market\_\allowbreak{}value\_\allowbreak{}musd\_\allowbreak{}gap & pregame\_\allowbreak{}available & outer\_\allowbreak{}fold\_\allowbreak{}train\_\allowbreak{}only;\allowbreak{} final=560 train ids \\
avg\_\allowbreak{}age\_\allowbreak{}mean & historical\_\allowbreak{}matches/\allowbreak{}teams/\allowbreak{}base\_\allowbreak{}predictions & avg\_\allowbreak{}age\_\allowbreak{}mean & pregame\_\allowbreak{}available & outer\_\allowbreak{}fold\_\allowbreak{}train\_\allowbreak{}only;\allowbreak{} final=560 train ids \\
avg\_\allowbreak{}age\_\allowbreak{}gap & historical\_\allowbreak{}matches/\allowbreak{}teams/\allowbreak{}base\_\allowbreak{}predictions & avg\_\allowbreak{}age\_\allowbreak{}gap & pregame\_\allowbreak{}available & outer\_\allowbreak{}fold\_\allowbreak{}train\_\allowbreak{}only;\allowbreak{} final=560 train ids \\
star\_\allowbreak{}index\_\allowbreak{}mean & historical\_\allowbreak{}matches/\allowbreak{}teams/\allowbreak{}base\_\allowbreak{}predictions & star\_\allowbreak{}index\_\allowbreak{}mean & pregame\_\allowbreak{}available & outer\_\allowbreak{}fold\_\allowbreak{}train\_\allowbreak{}only;\allowbreak{} final=560 train ids \\
star\_\allowbreak{}index\_\allowbreak{}gap & historical\_\allowbreak{}matches/\allowbreak{}teams/\allowbreak{}base\_\allowbreak{}predictions & star\_\allowbreak{}index\_\allowbreak{}gap & pregame\_\allowbreak{}available & outer\_\allowbreak{}fold\_\allowbreak{}train\_\allowbreak{}only;\allowbreak{} final=560 train ids \\
fan\_\allowbreak{}base\_\allowbreak{}index\_\allowbreak{}mean & historical\_\allowbreak{}matches/\allowbreak{}teams/\allowbreak{}base\_\allowbreak{}predictions & fan\_\allowbreak{}base\_\allowbreak{}index\_\allowbreak{}mean & pregame\_\allowbreak{}available & outer\_\allowbreak{}fold\_\allowbreak{}train\_\allowbreak{}only;\allowbreak{} final=560 train ids \\
fan\_\allowbreak{}base\_\allowbreak{}index\_\allowbreak{}gap & historical\_\allowbreak{}matches/\allowbreak{}teams/\allowbreak{}base\_\allowbreak{}predictions & fan\_\allowbreak{}base\_\allowbreak{}index\_\allowbreak{}gap & pregame\_\allowbreak{}available & outer\_\allowbreak{}fold\_\allowbreak{}train\_\allowbreak{}only;\allowbreak{} final=560 train ids \\
style\_\allowbreak{}attack\_\allowbreak{}mean & historical\_\allowbreak{}matches/\allowbreak{}teams/\allowbreak{}base\_\allowbreak{}predictions & style\_\allowbreak{}attack\_\allowbreak{}mean & pregame\_\allowbreak{}available & outer\_\allowbreak{}fold\_\allowbreak{}train\_\allowbreak{}only;\allowbreak{} final=560 train ids \\
style\_\allowbreak{}attack\_\allowbreak{}gap & historical\_\allowbreak{}matches/\allowbreak{}teams/\allowbreak{}base\_\allowbreak{}predictions & style\_\allowbreak{}attack\_\allowbreak{}gap & pregame\_\allowbreak{}available & outer\_\allowbreak{}fold\_\allowbreak{}train\_\allowbreak{}only;\allowbreak{} final=560 train ids \\
style\_\allowbreak{}defense\_\allowbreak{}mean & historical\_\allowbreak{}matches/\allowbreak{}teams/\allowbreak{}base\_\allowbreak{}predictions & style\_\allowbreak{}defense\_\allowbreak{}mean & pregame\_\allowbreak{}available & outer\_\allowbreak{}fold\_\allowbreak{}train\_\allowbreak{}only;\allowbreak{} final=560 train ids \\
style\_\allowbreak{}defense\_\allowbreak{}gap & historical\_\allowbreak{}matches/\allowbreak{}teams/\allowbreak{}base\_\allowbreak{}predictions & style\_\allowbreak{}defense\_\allowbreak{}gap & pregame\_\allowbreak{}available & outer\_\allowbreak{}fold\_\allowbreak{}train\_\allowbreak{}only;\allowbreak{} final=560 train ids \\
host\_\allowbreak{}flag\_\allowbreak{}mean & historical\_\allowbreak{}matches/\allowbreak{}teams/\allowbreak{}base\_\allowbreak{}predictions & host\_\allowbreak{}flag\_\allowbreak{}mean & pregame\_\allowbreak{}available & outer\_\allowbreak{}fold\_\allowbreak{}train\_\allowbreak{}only;\allowbreak{} final=560 train ids \\
host\_\allowbreak{}flag\_\allowbreak{}gap & historical\_\allowbreak{}matches/\allowbreak{}teams/\allowbreak{}base\_\allowbreak{}predictions & host\_\allowbreak{}flag\_\allowbreak{}gap & pregame\_\allowbreak{}available & outer\_\allowbreak{}fold\_\allowbreak{}train\_\allowbreak{}only;\allowbreak{} final=560 train ids \\
strength\_\allowbreak{}score\_\allowbreak{}mean & historical\_\allowbreak{}matches/\allowbreak{}teams/\allowbreak{}base\_\allowbreak{}predictions & strength\_\allowbreak{}score\_\allowbreak{}mean & pregame\_\allowbreak{}available & outer\_\allowbreak{}fold\_\allowbreak{}train\_\allowbreak{}only;\allowbreak{} final=560 train ids \\
strength\_\allowbreak{}score\_\allowbreak{}gap & historical\_\allowbreak{}matches/\allowbreak{}teams/\allowbreak{}base\_\allowbreak{}predictions & strength\_\allowbreak{}score\_\allowbreak{}gap & pregame\_\allowbreak{}available & outer\_\allowbreak{}fold\_\allowbreak{}train\_\allowbreak{}only;\allowbreak{} final=560 train ids \\
a\_\allowbreak{}hist\_\allowbreak{}goals\_\allowbreak{}roll5 & historical\_\allowbreak{}matches & a\_\allowbreak{}hist\_\allowbreak{}goals\_\allowbreak{}roll5 & strictly\_\allowbreak{}before\_\allowbreak{}kickoff\_\allowbreak{}shift1 & outer\_\allowbreak{}fold\_\allowbreak{}train\_\allowbreak{}only;\allowbreak{} final=560 train ids \\
a\_\allowbreak{}hist\_\allowbreak{}goals\_\allowbreak{}ewm5 & historical\_\allowbreak{}matches & a\_\allowbreak{}hist\_\allowbreak{}goals\_\allowbreak{}ewm5 & strictly\_\allowbreak{}before\_\allowbreak{}kickoff\_\allowbreak{}shift1 & outer\_\allowbreak{}fold\_\allowbreak{}train\_\allowbreak{}only;\allowbreak{} final=560 train ids \\
a\_\allowbreak{}hist\_\allowbreak{}xg\_\allowbreak{}roll5 & historical\_\allowbreak{}matches & a\_\allowbreak{}hist\_\allowbreak{}xg\_\allowbreak{}roll5 & strictly\_\allowbreak{}before\_\allowbreak{}kickoff\_\allowbreak{}shift1 & outer\_\allowbreak{}fold\_\allowbreak{}train\_\allowbreak{}only;\allowbreak{} final=560 train ids \\
a\_\allowbreak{}hist\_\allowbreak{}xg\_\allowbreak{}ewm5 & historical\_\allowbreak{}matches & a\_\allowbreak{}hist\_\allowbreak{}xg\_\allowbreak{}ewm5 & strictly\_\allowbreak{}before\_\allowbreak{}kickoff\_\allowbreak{}shift1 & outer\_\allowbreak{}fold\_\allowbreak{}train\_\allowbreak{}only;\allowbreak{} final=560 train ids \\
a\_\allowbreak{}hist\_\allowbreak{}shots\_\allowbreak{}roll5 & historical\_\allowbreak{}matches & a\_\allowbreak{}hist\_\allowbreak{}shots\_\allowbreak{}roll5 & strictly\_\allowbreak{}before\_\allowbreak{}kickoff\_\allowbreak{}shift1 & outer\_\allowbreak{}fold\_\allowbreak{}train\_\allowbreak{}only;\allowbreak{} final=560 train ids \\
a\_\allowbreak{}hist\_\allowbreak{}shots\_\allowbreak{}ewm5 & historical\_\allowbreak{}matches & a\_\allowbreak{}hist\_\allowbreak{}shots\_\allowbreak{}ewm5 & strictly\_\allowbreak{}before\_\allowbreak{}kickoff\_\allowbreak{}shift1 & outer\_\allowbreak{}fold\_\allowbreak{}train\_\allowbreak{}only;\allowbreak{} final=560 train ids \\
a\_\allowbreak{}hist\_\allowbreak{}possession\_\allowbreak{}roll5 & historical\_\allowbreak{}matches & a\_\allowbreak{}hist\_\allowbreak{}possession\_\allowbreak{}roll5 & strictly\_\allowbreak{}before\_\allowbreak{}kickoff\_\allowbreak{}shift1 & outer\_\allowbreak{}fold\_\allowbreak{}train\_\allowbreak{}only;\allowbreak{} final=560 train ids \\
a\_\allowbreak{}hist\_\allowbreak{}possession\_\allowbreak{}ewm5 & historical\_\allowbreak{}matches & a\_\allowbreak{}hist\_\allowbreak{}possession\_\allowbreak{}ewm5 & strictly\_\allowbreak{}before\_\allowbreak{}kickoff\_\allowbreak{}shift1 & outer\_\allowbreak{}fold\_\allowbreak{}train\_\allowbreak{}only;\allowbreak{} final=560 train ids \\
a\_\allowbreak{}hist\_\allowbreak{}attendance\_\allowbreak{}roll5 & historical\_\allowbreak{}matches & a\_\allowbreak{}hist\_\allowbreak{}attendance\_\allowbreak{}roll5 & strictly\_\allowbreak{}before\_\allowbreak{}kickoff\_\allowbreak{}shift1 & outer\_\allowbreak{}fold\_\allowbreak{}train\_\allowbreak{}only;\allowbreak{} final=560 train ids \\
a\_\allowbreak{}hist\_\allowbreak{}attendance\_\allowbreak{}ewm5 & historical\_\allowbreak{}matches & a\_\allowbreak{}hist\_\allowbreak{}attendance\_\allowbreak{}ewm5 & strictly\_\allowbreak{}before\_\allowbreak{}kickoff\_\allowbreak{}shift1 & outer\_\allowbreak{}fold\_\allowbreak{}train\_\allowbreak{}only;\allowbreak{} final=560 train ids \\
a\_\allowbreak{}hist\_\allowbreak{}tv\_\allowbreak{}viewers\_\allowbreak{}roll5 & historical\_\allowbreak{}matches & a\_\allowbreak{}hist\_\allowbreak{}tv\_\allowbreak{}viewers\_\allowbreak{}roll5 & strictly\_\allowbreak{}before\_\allowbreak{}kickoff\_\allowbreak{}shift1 & outer\_\allowbreak{}fold\_\allowbreak{}train\_\allowbreak{}only;\allowbreak{} final=560 train ids \\
a\_\allowbreak{}hist\_\allowbreak{}tv\_\allowbreak{}viewers\_\allowbreak{}ewm5 & historical\_\allowbreak{}matches & a\_\allowbreak{}hist\_\allowbreak{}tv\_\allowbreak{}viewers\_\allowbreak{}ewm5 & strictly\_\allowbreak{}before\_\allowbreak{}kickoff\_\allowbreak{}shift1 & outer\_\allowbreak{}fold\_\allowbreak{}train\_\allowbreak{}only;\allowbreak{} final=560 train ids \\
b\_\allowbreak{}hist\_\allowbreak{}goals\_\allowbreak{}roll5 & historical\_\allowbreak{}matches & b\_\allowbreak{}hist\_\allowbreak{}goals\_\allowbreak{}roll5 & strictly\_\allowbreak{}before\_\allowbreak{}kickoff\_\allowbreak{}shift1 & outer\_\allowbreak{}fold\_\allowbreak{}train\_\allowbreak{}only;\allowbreak{} final=560 train ids \\
b\_\allowbreak{}hist\_\allowbreak{}goals\_\allowbreak{}ewm5 & historical\_\allowbreak{}matches & b\_\allowbreak{}hist\_\allowbreak{}goals\_\allowbreak{}ewm5 & strictly\_\allowbreak{}before\_\allowbreak{}kickoff\_\allowbreak{}shift1 & outer\_\allowbreak{}fold\_\allowbreak{}train\_\allowbreak{}only;\allowbreak{} final=560 train ids \\
b\_\allowbreak{}hist\_\allowbreak{}xg\_\allowbreak{}roll5 & historical\_\allowbreak{}matches & b\_\allowbreak{}hist\_\allowbreak{}xg\_\allowbreak{}roll5 & strictly\_\allowbreak{}before\_\allowbreak{}kickoff\_\allowbreak{}shift1 & outer\_\allowbreak{}fold\_\allowbreak{}train\_\allowbreak{}only;\allowbreak{} final=560 train ids \\
b\_\allowbreak{}hist\_\allowbreak{}xg\_\allowbreak{}ewm5 & historical\_\allowbreak{}matches & b\_\allowbreak{}hist\_\allowbreak{}xg\_\allowbreak{}ewm5 & strictly\_\allowbreak{}before\_\allowbreak{}kickoff\_\allowbreak{}shift1 & outer\_\allowbreak{}fold\_\allowbreak{}train\_\allowbreak{}only;\allowbreak{} final=560 train ids \\
b\_\allowbreak{}hist\_\allowbreak{}shots\_\allowbreak{}roll5 & historical\_\allowbreak{}matches & b\_\allowbreak{}hist\_\allowbreak{}shots\_\allowbreak{}roll5 & strictly\_\allowbreak{}before\_\allowbreak{}kickoff\_\allowbreak{}shift1 & outer\_\allowbreak{}fold\_\allowbreak{}train\_\allowbreak{}only;\allowbreak{} final=560 train ids \\
b\_\allowbreak{}hist\_\allowbreak{}shots\_\allowbreak{}ewm5 & historical\_\allowbreak{}matches & b\_\allowbreak{}hist\_\allowbreak{}shots\_\allowbreak{}ewm5 & strictly\_\allowbreak{}before\_\allowbreak{}kickoff\_\allowbreak{}shift1 & outer\_\allowbreak{}fold\_\allowbreak{}train\_\allowbreak{}only;\allowbreak{} final=560 train ids \\
b\_\allowbreak{}hist\_\allowbreak{}possession\_\allowbreak{}roll5 & historical\_\allowbreak{}matches & b\_\allowbreak{}hist\_\allowbreak{}possession\_\allowbreak{}roll5 & strictly\_\allowbreak{}before\_\allowbreak{}kickoff\_\allowbreak{}shift1 & outer\_\allowbreak{}fold\_\allowbreak{}train\_\allowbreak{}only;\allowbreak{} final=560 train ids \\
b\_\allowbreak{}hist\_\allowbreak{}possession\_\allowbreak{}ewm5 & historical\_\allowbreak{}matches & b\_\allowbreak{}hist\_\allowbreak{}possession\_\allowbreak{}ewm5 & strictly\_\allowbreak{}before\_\allowbreak{}kickoff\_\allowbreak{}shift1 & outer\_\allowbreak{}fold\_\allowbreak{}train\_\allowbreak{}only;\allowbreak{} final=560 train ids \\
b\_\allowbreak{}hist\_\allowbreak{}attendance\_\allowbreak{}roll5 & historical\_\allowbreak{}matches & b\_\allowbreak{}hist\_\allowbreak{}attendance\_\allowbreak{}roll5 & strictly\_\allowbreak{}before\_\allowbreak{}kickoff\_\allowbreak{}shift1 & outer\_\allowbreak{}fold\_\allowbreak{}train\_\allowbreak{}only;\allowbreak{} final=560 train ids \\
b\_\allowbreak{}hist\_\allowbreak{}attendance\_\allowbreak{}ewm5 & historical\_\allowbreak{}matches & b\_\allowbreak{}hist\_\allowbreak{}attendance\_\allowbreak{}ewm5 & strictly\_\allowbreak{}before\_\allowbreak{}kickoff\_\allowbreak{}shift1 & outer\_\allowbreak{}fold\_\allowbreak{}train\_\allowbreak{}only;\allowbreak{} final=560 train ids \\
b\_\allowbreak{}hist\_\allowbreak{}tv\_\allowbreak{}viewers\_\allowbreak{}roll5 & historical\_\allowbreak{}matches & b\_\allowbreak{}hist\_\allowbreak{}tv\_\allowbreak{}viewers\_\allowbreak{}roll5 & strictly\_\allowbreak{}before\_\allowbreak{}kickoff\_\allowbreak{}shift1 & outer\_\allowbreak{}fold\_\allowbreak{}train\_\allowbreak{}only;\allowbreak{} final=560 train ids \\
b\_\allowbreak{}hist\_\allowbreak{}tv\_\allowbreak{}viewers\_\allowbreak{}ewm5 & historical\_\allowbreak{}matches & b\_\allowbreak{}hist\_\allowbreak{}tv\_\allowbreak{}viewers\_\allowbreak{}ewm5 & strictly\_\allowbreak{}before\_\allowbreak{}kickoff\_\allowbreak{}shift1 & outer\_\allowbreak{}fold\_\allowbreak{}train\_\allowbreak{}only;\allowbreak{} final=560 train ids \\
worldcup\_\allowbreak{}stage & historical\_\allowbreak{}matches/\allowbreak{}teams/\allowbreak{}base\_\allowbreak{}predictions & worldcup\_\allowbreak{}stage & pregame\_\allowbreak{}available & outer\_\allowbreak{}fold\_\allowbreak{}train\_\allowbreak{}only;\allowbreak{} final=560 train ids \\
fan\_\allowbreak{}stage\_\allowbreak{}interaction & historical\_\allowbreak{}matches/\allowbreak{}teams/\allowbreak{}base\_\allowbreak{}predictions & fan\_\allowbreak{}stage\_\allowbreak{}interaction & pregame\_\allowbreak{}available & outer\_\allowbreak{}fold\_\allowbreak{}train\_\allowbreak{}only;\allowbreak{} final=560 train ids \\
star\_\allowbreak{}stage\_\allowbreak{}interaction & historical\_\allowbreak{}matches/\allowbreak{}teams/\allowbreak{}base\_\allowbreak{}predictions & star\_\allowbreak{}stage\_\allowbreak{}interaction & pregame\_\allowbreak{}available & outer\_\allowbreak{}fold\_\allowbreak{}train\_\allowbreak{}only;\allowbreak{} final=560 train ids \\
strength\_\allowbreak{}entropy\_\allowbreak{}interaction & historical\_\allowbreak{}matches/\allowbreak{}teams/\allowbreak{}base\_\allowbreak{}predictions & strength\_\allowbreak{}entropy\_\allowbreak{}interaction & pregame\_\allowbreak{}available & outer\_\allowbreak{}fold\_\allowbreak{}train\_\allowbreak{}only;\allowbreak{} final=560 train ids \\
\end{longtable}
\par
}


\begingroup
\let\oldsmall\small
\renewcommand{\small}{\scriptsize}
\renewcommand{\arraystretch}{0.78}
\begin{table}[H]
\centering
\caption{问题一时序交叉验证误差}
\label{tab:q1-model-metrics}
\small
\resizebox{\textwidth}{!}{%
\begin{tabular}{llrrrrr}
\toprule
模型 & 窗口 & MSE & RMSE & MAE & 训练截至 & 验证起始 \\
\midrule
MeanBaseline & 折1 & 598.919 & 24.473 & 19.751 & 2019-01-26 & 2019-02-03 \\
ElasticNet & 折1 & 243.138 & 15.593 & 12.757 & 2019-01-26 & 2019-02-03 \\
ExtraTrees & 折1 & 298.141 & 17.267 & 14.214 & 2019-01-26 & 2019-02-03 \\
CatBoost & 折1 & 274.080 & 16.555 & 13.172 & 2019-01-26 & 2019-02-03 \\
MeanBaseline & 折2 & 621.856 & 24.937 & 19.817 & 2020-03-29 & 2020-04-04 \\
ElasticNet & 折2 & 226.619 & 15.054 & 11.929 & 2020-03-29 & 2020-04-04 \\
ExtraTrees & 折2 & 261.458 & 16.170 & 12.694 & 2020-03-29 & 2020-04-04 \\
CatBoost & 折2 & 244.634 & 15.641 & 12.265 & 2020-03-29 & 2020-04-04 \\
MeanBaseline & 折3 & 528.105 & 22.981 & 18.864 & 2021-04-19 & 2021-04-22 \\
ElasticNet & 折3 & 239.340 & 15.471 & 12.203 & 2021-04-19 & 2021-04-22 \\
ExtraTrees & 折3 & 256.627 & 16.020 & 13.012 & 2021-04-19 & 2021-04-22 \\
CatBoost & 折3 & 243.231 & 15.596 & 12.693 & 2021-04-19 & 2021-04-22 \\
MeanBaseline & 折4 & 575.265 & 23.985 & 18.588 & 2022-07-03 & 2022-07-05 \\
ElasticNet & 折4 & 214.684 & 14.652 & 11.510 & 2022-07-03 & 2022-07-05 \\
ExtraTrees & 折4 & 242.182 & 15.562 & 12.326 & 2022-07-03 & 2022-07-05 \\
CatBoost & 折4 & 230.889 & 15.195 & 11.874 & 2022-07-03 & 2022-07-05 \\
MeanBaseline & 折5 & 581.872 & 24.122 & 19.646 & 2023-07-28 & 2023-07-29 \\
ElasticNet & 折5 & 176.646 & 13.291 & 10.739 & 2023-07-28 & 2023-07-29 \\
ExtraTrees & 折5 & 192.403 & 13.871 & 11.426 & 2023-07-28 & 2023-07-29 \\
CatBoost & 折5 & 178.211 & 13.350 & 10.915 & 2023-07-28 & 2023-07-29 \\
ElasticNet & 总体均值 & 220.086 & 14.812 & 11.828 & -- & -- \\
CatBoost & 总体均值 & 234.209 & 15.267 & 12.184 & -- & -- \\
ExtraTrees & 总体均值 & 250.162 & 15.778 & 12.734 & -- & -- \\
MeanBaseline & 总体均值 & 581.203 & 24.099 & 19.333 & -- & -- \\
\bottomrule
\end{tabular}
}
\vspace{2pt}
\noindent\begin{minipage}{0.96\textwidth}\footnotesize 注：误差单位为百万人口径；最终选择模型为 ElasticNet。\end{minipage}
\end{table}

\endgroup

\begin{table}[H]
\centering
\caption{问题二场馆负载与赛程汇总}
\label{tab:q2-schedule-summary}
\small
\resizebox{\textwidth}{!}{%
\begin{tabular}{llrrrr}
\toprule
场馆 & 城市 & 场次 & 黄金时段场次 & 平均旅行指数 & 预计观众合计（万人） \\
\midrule
V01 & Atlanta & 8 & 4 & 0.254 & 32.6 \\
V02 & Boston & 8 & 5 & 0.321 & 30.0 \\
V03 & Dallas & 7 & 7 & 0.081 & 27.1 \\
V04 & Guadalajara & 8 & 6 & 0.152 & 35.1 \\
V05 & Houston & 8 & 6 & 0.086 & 35.2 \\
V06 & Kansas City & 5 & 5 & 0.086 & 17.3 \\
V07 & Los Angeles & 4 & 1 & 0.270 & 16.7 \\
V08 & Mexico City & 7 & 7 & 0.110 & 34.3 \\
V09 & Miami & 3 & 1 & 0.244 & 13.9 \\
V10 & Monterrey & 2 & 0 & 0.096 & 6.7 \\
V11 & New York/New Jersey & 2 & 1 & 0.117 & 6.9 \\
V12 & Philadelphia & 2 & 1 & 0.272 & 9.3 \\
V13 & San Francisco Bay Area & 2 & 1 & 0.414 & 9.0 \\
V14 & Seattle & 2 & 0 & 0.556 & 8.4 \\
V15 & Toronto & 2 & 0 & 0.266 & 8.5 \\
V16 & Vancouver & 2 & 0 & 0.504 & 9.4 \\
\bottomrule
\end{tabular}
}
\vspace{2pt}
\noindent\begin{minipage}{0.96\textwidth}\footnotesize 注：主方案 Z2=0.484946，最小轮间休息 63 小时。\end{minipage}
\end{table}


\begin{table}[H]
\centering
\caption{问题三晋级概率、条件概率与风险}
\label{tab:q3-probability-risk}
\small
\resizebox{\textwidth}{!}{%
\begin{tabular}{llrlrrrrr}
\toprule
比赛 & A队 & A队晋级 & B队 & B队晋级 & A胜条件晋级 & B胜条件晋级 & 无悬念风险 & 默契风险 \\
\midrule
GA05 & Mexico & 1.000 & Czechia & 1.000 & 1.000 & 1.000 & 1.000 & 0.650 \\
GA06 & South Africa & 0.226 & South Korea & 0.541 & 1.000 & 0.991 & 0.154 & 0.000 \\
GB05 & Canada & 0.355 & Switzerland & 1.000 & 1.000 & 1.000 & 0.542 & 0.803 \\
GB06 & Bosnia and Herzegovina & 0.601 & Qatar & 0.318 & 1.000 & 1.000 & 0.087 & 0.013 \\
GC05 & Brazil & 1.000 & Scotland & 0.077 & 1.000 & 1.000 & 0.857 & 0.015 \\
GC06 & Morocco & 0.963 & Haiti & 0.209 & 1.000 & 1.000 & 0.599 & 0.704 \\
GD05 & United States & 1.000 & Turkey & 1.000 & 1.000 & 1.000 & 1.000 & 0.650 \\
GD06 & Paraguay & 0.314 & Australia & 0.631 & 0.851 & 1.000 & 0.103 & 0.000 \\
GE05 & Germany & 1.000 & Ecuador & 0.954 & 1.000 & 1.000 & 0.912 & 0.681 \\
GE06 & Curacao & 0.002 & Ivory Coast & 1.000 & 0.021 & 1.000 & 0.995 & 0.000 \\
GF05 & Netherlands & 1.000 & Tunisia & 0.070 & 1.000 & 1.000 & 0.870 & 0.002 \\
GF06 & Japan & 0.934 & Sweden & 0.319 & 1.000 & 1.000 & 0.442 & 0.001 \\
GG05 & Belgium & 1.000 & New Zealand & 0.782 & 1.000 & 1.000 & 0.659 & 0.769 \\
GG06 & Egypt & 0.808 & Iran & 0.155 & 1.000 & 0.416 & 0.428 & 0.000 \\
GH05 & Spain & 1.000 & Uruguay & 0.986 & 1.000 & 1.000 & 0.972 & 0.660 \\
GH06 & Cape Verde & 0.025 & Saudi Arabia & 0.965 & 0.100 & 1.000 & 0.885 & 0.000 \\
GI05 & France & 1.000 & Norway & 1.000 & 1.000 & 1.000 & 1.000 & 0.650 \\
GI06 & Senegal & 0.886 & Iraq & 0.056 & 1.000 & 0.496 & 0.692 & 0.000 \\
GJ05 & Argentina & 1.000 & Jordan & 1.000 & 1.000 & 1.000 & 1.000 & 0.650 \\
GJ06 & Algeria & 0.220 & Austria & 0.510 & 0.986 & 0.995 & 0.157 & 0.000 \\
GK05 & Portugal & 1.000 & Colombia & 1.000 & 1.000 & 1.000 & 1.000 & 0.650 \\
GK06 & DR Congo & 0.288 & Uzbekistan & 0.122 & 0.956 & 0.292 & 0.375 & 0.000 \\
GL05 & England & 1.000 & Panama & 0.588 & 1.000 & 1.000 & 0.516 & 0.820 \\
GL06 & Croatia & 1.000 & Ghana & 0.093 & 1.000 & 0.614 & 0.832 & 0.000 \\
\bottomrule
\end{tabular}
}
\vspace{2pt}
\noindent\begin{minipage}{0.96\textwidth}\footnotesize 注：条件晋级概率来自同一 20000 次联合蒙特卡洛样本，晋级名额质量守恒为每次 32 队。\end{minipage}
\end{table}


{
\footnotesize
\hfuzz=0.2pt
\setlength{\tabcolsep}{3pt}
\begin{longtable}{@{}p{0.08\textwidth}p{0.26\textwidth}p{0.08\textwidth}p{0.11\textwidth}p{0.18\textwidth}p{0.209\textwidth}@{}}
\caption{问题四评价指标定义与可比口径}\label{tab:q4-metric-definition}\\
\toprule
指标 & 定义 & 单位 & 偏好方向 & 数据源 & 口径说明 \\
\midrule
\endfirsthead
\multicolumn{6}{c}{续表} \\
\toprule
指标 & 定义 & 单位 & 偏好方向 & 数据源 & 口径说明 \\
\midrule
\endhead
\midrule \multicolumn{6}{r}{续下页} \\
\endfoot
\bottomrule
\endlastfoot
rest\_\allowbreak{}min\_\allowbreak{}hours & 球队相邻比赛 UTC 间隔的最小值 & 小时 & 越高越优 & 统一结构评价函数 & 现实赛程直接复算 \\
rest\_\allowbreak{}mean\_\allowbreak{}hours & 球队相邻比赛 UTC 间隔的平均值 & 小时 & 越高越优 & 统一结构评价函数 & 现实赛程直接复算 \\
timezone\_\allowbreak{}crossings & 相邻比赛场馆 UTC 偏移发生变化的次数 & 次 & 越低越优 & 统一结构评价函数 & 现实赛程直接复算 \\
timezone\_\allowbreak{}crossing\_\allowbreak{}rate & 跨时区次数除以球队转场次数 & 比例 & 越低越优 & 统一结构评价函数 & 现实赛程直接复算 \\
travel\_\allowbreak{}mean\_\allowbreak{}km & 相邻比赛实际场馆大圆距离的平均值 & 公里/次 & 越低越优 & 统一结构评价函数 & 现实赛程直接复算 \\
venue\_\allowbreak{}change\_\allowbreak{}rate & 相邻轮更换场馆的球队转场比例 & 比例 & 越低越优 & 统一结构评价函数 & 现实赛程直接复算 \\
venue\_\allowbreak{}utilization\_\allowbreak{}rate & 发生比赛的场馆日占可用场馆日比例 & 比例 & 越高越优 & 统一结构评价函数 & 现实赛程直接复算 \\
venue\_\allowbreak{}daily\_\allowbreak{}peak & 单一场馆单日承办场次最大值 & 场/场馆日 & 越低越优 & 统一结构评价函数 & 现实赛程直接复算 \\
capacity\_\allowbreak{}occupancy\_\allowbreak{}rate & 预计现场观众与场馆容量之比的场均值 & 比例 & 越高越优 & 统一结构评价函数 & 需求端为题内代理 \\
prime\_\allowbreak{}coverage\_\allowbreak{}rate & 按赛事进度日和 UTC 钟点映射后落入前四分位时段的比例 & 比例 & 越高越优 & 统一结构评价函数 & 赛事进度与UTC时刻映射 \\
expected\_\allowbreak{}attendance\_\allowbreak{}per\_\allowbreak{}match & 预计现场观众总量除以比赛场次 & 人/场 & 越高越优 & 统一结构评价函数 & 需求端为题内代理 \\
T & $T=\mathrm{norm}(\sum \mathrm{ticket})$ & 无量纲 & 越高越优 & 统一代理评价函数 & 共同包络Min--Max \\
B & $B=\mathrm{norm}(\sum \mathrm{broadcast})$ & 无量纲 & 越高越优 & 统一代理评价函数 & 共同包络Min--Max \\
U & $U=\overline{\mathrm{uncertainty}}$ & 无量纲 & 越高越优 & 统一代理评价函数 & 题内代理评价 \\
H & $H=\overline{\mathrm{attractiveness}}/100$ & 无量纲 & 越高越优 & 统一代理评价函数 & 题内代理评价 \\
C & $C=\mathrm{norm}(\mathrm{setup+operation+security})$ & 无量纲 & 越低越优 & 统一代理评价函数 & 共同包络Min--Max \\
D & $D=\overline{\mathrm{travel}}$ & 无量纲 & 越低越优 & 统一代理评价函数 & 题内代理评价 \\
F & $F=(\Delta prime+\Delta large)/6$ & 无量纲 & 越低越优 & 统一代理评价函数 & 题内代理评价 \\
R & $R=\overline{\mathrm{risk}}$ & 无量纲 & 越低越优 & 统一代理评价函数 & 题内代理评价 \\
Z4 & $Z4=0.25T+0.25B+0.15U+0.10H-0.08C-0.07D-0.06F-0.04R$ & 无量纲 & 越高越优 & 统一代理评价函数 & 条件性加权总分 \\
\end{longtable}
\noindent\begin{minipage}{0.96\textwidth}\footnotesize 注：结构层报告物理量；代理层用完整赛前Elo映射比赛、仅按容量和安保可行域映射场馆，黄金时段按赛事进度日与UTC钟点唯一映射；T/B/C在全部报告场景共用的包络边界上归一化。\end{minipage}
\par
}


\subsection{完整可运行源程序}

下列清单给出模型与正式复核链所使用的完整源程序，不省略函数或执行入口。图表生成程序的完整文件同时保存在支撑材料中。

\lstinputlisting[caption={一键复现入口},basicstyle=\tiny\ttfamily]{../reproduce.py}
\lstinputlisting[caption={四问总入口},basicstyle=\tiny\ttfamily]{../code/main.py}
\lstinputlisting[caption={官方附件摄入检查},basicstyle=\tiny\ttfamily]{../code/data_check.py}
\lstinputlisting[caption={统一参数定义},basicstyle=\tiny\ttfamily]{../code/params.py}
\lstinputlisting[caption={共享数据与运行工具},basicstyle=\tiny\ttfamily]{../code/utils.py}
\lstinputlisting[caption={问题一完整程序},basicstyle=\tiny\ttfamily]{../code/problem1.py}
\lstinputlisting[caption={问题二完整程序},basicstyle=\tiny\ttfamily]{../code/problem2.py}
\lstinputlisting[caption={问题三完整程序},basicstyle=\tiny\ttfamily]{../code/problem3.py}
\lstinputlisting[caption={问题四完整程序},basicstyle=\tiny\ttfamily]{../code/problem4.py}
\lstinputlisting[caption={灵敏度分析完整程序},basicstyle=\tiny\ttfamily]{../code/sensitivity_analysis.py}
\lstinputlisting[caption={正式输出约束复核程序},basicstyle=\tiny\ttfamily]{../code/constraint_audit.py}
\lstinputlisting[caption={图表统一生成入口},basicstyle=\tiny\ttfamily]{../figures/generate_all.py}
\lstinputlisting[caption={图表数据整理程序},basicstyle=\tiny\ttfamily]{../figures/prep_plot_data.py}
\lstinputlisting[caption={统一绘图实现},basicstyle=\tiny\ttfamily]{../figures/plot_common.py}
\lstinputlisting[caption={论文表格生成程序},basicstyle=\tiny\ttfamily]{../figures/generate_tables.py}
